\chapter{A Map of the Industry}\label{ch:in:a-map-of-the-industry}

The United States government's industry classification has an entry for ``market making for securities''. It files it
under code 523150, Investment Banking and Securities Intermediation, beside the underwriting departments of the largest
banks. A firm that trades only its own money is a ``securities speculator for own account'' and goes to 523910,
Miscellaneous Intermediation, with ``individuals investing in financial contracts on own account''. On the other side of
the same securities industry, the SEC's code 6211 covers an investment bank, an electronic market maker and an asset
manager with \$14 trillion under management alike. None of these codes says whose capital is at risk, who pays the firm, or how long
it holds a position, and those three answers decide what a job there is. This chapter builds the map the rest of the
book uses, and shows how to place an employer on it from public records alone.

\section{Five questions that place an employer}

Book~1 defined the firms that trade: the market maker, the proprietary trading firm, the hedge fund and the
multi-manager platform, the asset manager, and the bank's dealing desk. For someone choosing where to work, five
questions separate them more finely than any name.

\begin{enumerate}
\item \textbf{Whose capital is at risk?} The owners' (a proprietary firm), outside investors' (a fund), clients' under
a mandate (an asset manager), shareholders' within a regulated balance sheet (a bank).
\item \textbf{Who pays the firm?} Counterparties through the spread, investors through fees on assets and on profits,
clients through commissions, or nobody but the market.
\item \textbf{How long is a position held?} Seconds for a market maker, minutes to days for a medium-frequency firm,
months for a trend follower, years for an asset owner.
\item \textbf{What does it sell, and in which markets?} A service (liquidity, execution, a product, a licence) in one or
several asset classes.
\item \textbf{How large is it?} In people, in capital and in revenue: tens, hundreds, thousands or hundreds of thousands
of employees.
\end{enumerate}

\begin{definition}[Employer taxonomy]\label{def:in:a-map-of-the-industry:taxonomy}
An \emph{employer taxonomy}\index{employer taxonomy} of the markets industry classifies an employer, or one business line of
a group, by its business model -- the answers to the first three questions -- and then by product and size. This book
uses thirteen business models in four groups: principal trading (market maker, medium-frequency proprietary firm, crypto
trading firm), investment manager (systematic fund, multi-manager platform, discretionary fund, asset manager, asset
owner), bank, and infrastructure (exchange, broker, vendor, regulator).
\end{definition}

\begin{center}\small
\begin{tabular}{@{}llp{4.2cm}l@{}}
\toprule
business model & whose capital & who pays & holding period\\
\midrule
market maker & owners' & the spread and rebates & seconds to hours\\
medium-frequency proprietary firm & owners' & trading profit & minutes to days\\
systematic fund & investors' & fees on assets and profits & minutes to months\\
multi-manager platform & investors' & pass-through costs and profits & days to months\\
asset manager & clients' & fees on assets & months to years\\
asset owner & beneficiaries' & its own budget & years\\
bank markets division & shareholders' & spread and client business & minutes to months\\
exchange, broker, vendor & shareholders' & transaction, data and licence fees & none\\
\bottomrule
\end{tabular}
\end{center}

\begin{example}[Three employers, five answers]\label{ex:in:a-map-of-the-industry:three}
A commodity trading advisor running trend-following futures programmes risks investors' capital, is paid a management
and a performance fee, holds positions for weeks to months, sells a managed-futures product in every liquid futures
market, and may employ a few dozen people. An exchange-traded-fund market maker risks its owners' capital, is paid by
the spread between the fund's price and its basket, holds positions for minutes to hours, and sells liquidity in the
funds it quotes. A bank's rates flow desk risks shareholders' capital inside a regulated balance sheet, is paid through
the spread and the client relationship, holds positions for days, and sells prices to institutional clients. Three
jobs called ``trader'' with nothing in common but the screens.
\end{example}

\begin{remark}[Groups and legal entities]\label{rem:in:a-map-of-the-industry:groups}
A group may own a market maker, a fund and a technology subsidiary under one brand, and a bank's markets division
contains businesses of several kinds. The taxonomy classifies the legal entity or the business line an employee joins,
not the brand. Where a group's businesses differ, its job offers do too.
\end{remark}

\section{The taxonomy: business models and their sub-types}

\Cref{fig:in:a-map-of-the-industry:tree} draws the thirteen models in their four groups, with the chapters of this book
that describe each as an employer. The principal-trading group is where the exchange-industry associations draw their
own line.

\begin{definition}[Principal trading firm]\label{def:in:a-map-of-the-industry:ptf}
A \emph{principal trading firm}\index{principal trading firm} is a proprietary trading firm (One Quant Book~1) in the sense
the exchange-industry associations use: a firm that trades only its own capital, on exchanges and electronic venues,
mostly as a provider of liquidity, and has no clients whose money it invests.
\end{definition}

The United States association limits membership of its principal traders' group to ``firms trading their own
capital''; the European association describes its members as ``independent market makers and providers of liquidity
and risk transfer''. The term matters to an employee because it fixes two things at once: every dollar of revenue is
trading revenue, and nobody but the owners has a claim on the profit (Book~16, chapter~2).

\begin{omfigure}
\begin{tikzpicture}[font=\footnotesize,>=stealth,
  grp/.style={draw,rounded corners,fill=omDef!12,minimum width=2.9cm,minimum height=0.6cm,align=center},
  mdl/.style={draw,rounded corners,fill=omProp!8,minimum width=5.9cm,minimum height=0.42cm,align=center,inner sep=2pt}]
\foreach \m/\y/\c in {a/0/{market maker (ch.~2, 3)},b/-0.55/{medium-frequency proprietary firm (ch.~6)},c/-1.1/{crypto trading firm (ch.~10)},
  d/-1.9/{systematic fund (ch.~4)},e/-2.45/{multi-manager platform (ch.~5)},f/-3.0/{discretionary fund (ch.~4, 5)},
  g/-3.55/{asset manager (ch.~8)},h/-4.1/{asset owner (ch.~8)},i/-4.9/{bank markets division (ch.~7)},
  j/-5.7/{exchange (ch.~9)},k/-6.25/{broker (ch.~9)},l/-6.8/{vendor (ch.~9)},n/-7.35/{regulator (ch.~9)}}
  {\node[mdl] (\m) at (6.2,\y) {\c};}
\node[grp] (P) at (0,-0.55) {principal trading};
\node[grp] (I) at (0,-3.0) {investment manager};
\node[grp] (B) at (0,-4.9) {bank};
\node[grp] (N) at (0,-6.525) {infrastructure};
\foreach \g/\ms in {P/{a,b,c},I/{d,e,f,g,h},B/{i},N/{j,k,l,n}} {\foreach \m in \ms {\draw[->,gray] (\g.east) -- (\m.west);}}
\node[font=\scriptsize,align=center] at (0,0.55) {group};
\node[font=\scriptsize,align=center] at (6.2,0.55) {business model (where this book describes it)};
\end{tikzpicture}

\omcaption{The \omterm{def:in:a-map-of-the-industry:taxonomy}{employer taxonomy} of this book: thirteen business models in four groups, and the chapters of Part~I that
describe each as an employer. The roles that cut across them are Part~III's.}\label{fig:in:a-map-of-the-industry:tree}
\end{omfigure}

Within each model, sub-types differ enough to change the job. Among market makers, an electronic delta-one firm quoting
shares, futures and currencies (chapter~2) is not an options house pricing thousands of strikes (chapter~3). Among
investment managers, a research-driven multi-strategy fund, a trend follower and a statistical-arbitrage house share a
method but not a pace (chapter~4). Banks differ by product strength: some earn most of their markets revenue in equity
derivatives, most in rates and credit (chapter~7). The map is also a map of this series. The table below lists, for each business model, the books a newcomer to it should read.

\begin{center}\small
\begin{tabular}{@{}lp{7.2cm}@{}}
\toprule
business model & One Quant Books that teach its work\\
\midrule
market maker & 1 (markets), 10 (microstructure), 11 (market making), 13--14 (systems, networks)\\
medium-frequency proprietary firm & 7 (research craft), 8 (strategies), 12 (machine learning)\\
crypto trading firm & 3 (crypto markets), 11, 14 (crypto connectivity)\\
systematic fund & 4 (methods), 7, 8, 9 (strategies), 12\\
multi-manager platform & 8, 16 (the desk and the firm)\\
asset manager and owner & 1, 8 (the asset managers' strategies)\\
bank markets division & 2 (rates, FX, credit), 5 (derivatives), 6 (XVA and risk), 9 (bank desks)\\
exchange, broker, vendor, regulator & 1, 10, 14, 15 (platforms), 16\\
\bottomrule
\end{tabular}
\end{center}

\section{Size: headcount, capital and revenue}

Size is the fifth question and the most visible one. It ranges over five orders of magnitude, from trading firms of a
dozen people to banks of hundreds of thousands, and it changes the job more than it seems: a firm of fifty has no
separate risk department, one of fifty thousand has several.

\begin{dated}{2026-09}{How large employers are}\label{dat:in:a-map-of-the-industry:sizes}
\textbf{Headcounts from filings and firms' own pages} (end of 2025 unless stated): JPMorganChase 318\,512 employees;
Goldman Sachs a headcount of 47\,400; BlackRock about 24\,900; CME Group about 3\,875; Virtu Financial about 1\,027
(February 2026). Two multi-manager platforms state on their home pages ``7\,000+ employees globally'' and ``360+
investment teams'' (Millennium) and ``2000+ investment team and support professionals'' (Balyasny).
\textbf{Counts of firms} (FINRA, end of 2024): 3\,249 registered broker-dealers, of which 2\,891 have 1--150 registered
representatives, 209 have 151--499 and 149 have 500 or more; 32\,090 firms registered only as investment advisers.
FINRA places 103 broker-dealers in its proprietary-trading and market-making business segments, 93 of them small by
its definition and none large.
\end{dated}

Two lessons follow. The number of registered investment advisers is ten times the number of broker-dealers, and it
grew while the number of broker-dealers fell (\cref{fig:in:a-map-of-the-industry:reg}): most firms in the industry
manage money, and most are small. And the \omterm{def:in:a-map-of-the-industry:ptf}{principal trading firms} are few and mostly small by the regulator's count,
which measures registered representatives, not staff: a market maker with a thousand engineers and traders may register
only the few who deal with the public, so its FINRA size says nothing about its real size. Every size has a unit, and
the unit has to be read.

\begin{omfigure}
\begin{tikzpicture}
\begin{axis}[width=11cm,height=5.4cm,xlabel={\small year end},ylabel={\small firms (thousands)},xmin=2014.6,xmax=2024.4,
  ymin=0,ymax=36,xtick={2015,2016,2017,2018,2019,2020,2021,2022,2023,2024},xticklabel style={/pgf/number format/1000 sep={}},
  tick label style={font=\footnotesize},grid=major,grid style={gray!20},table/col sep=comma,
  legend style={legend cell align=left,font=\scriptsize,at={(0.5,-0.3)},anchor=north,/tikz/every even column/.append style={column sep=8pt}},legend columns=2]
\addplot[omDef,thick,mark=*] table[x=year,y=ia_only]{figdata/industry/01-a-map-of-the-industry/registrations.csv};
\addplot[omThm,thick,mark=square*] table[x=year,y=bd]{figdata/industry/01-a-map-of-the-industry/registrations.csv};
\legend{registered only as investment advisers,FINRA-registered broker-dealers}
\end{axis}
\end{tikzpicture}

\omcaption{Firms registered in the US securities industry, 2015--2024: broker-dealers (including those also registered
as advisers) fell from 3.94 to 3.25 thousand, firms registered only as investment advisers rose from 28.7 to 32.1
thousand. Data: FINRA Industry Snapshot 2025, table 2.1.5, through \texttt{in\_map.finra}.}\label{fig:in:a-map-of-the-industry:reg}
\end{omfigure}

\begin{remark}[Size bands in this book]\label{rem:in:a-map-of-the-industry:bands}
Where this book sorts employers by size it uses staff, not registered representatives, in four bands: under 50, 50 to
500, 500 to 5\,000, and above 5\,000. A headcount is a dated range, never a point: firms round, count contractors
differently, and change size faster than their filings do.
\end{remark}

\section{Reading an employer from outside}

An outsider can place an employer from three kinds of public record: the industry codes that statistical agencies and
regulators assign, the registrations and periodic reports the law requires, and the firm's own publications.

\begin{definition}[Industry classification code]\label{def:in:a-map-of-the-industry:code}
An \emph{industry classification code}\index{industry classification code} assigns an establishment or a company to one
industry of a statistical classification by its main activity: the North American Industry Classification System
(NAICS, six digits), the Standard Industrial Classification that the SEC still uses on EDGAR (SIC, four digits), or the
European Union's NACE.
\end{definition}

The codes are coarse where the industry is fine. NAICS 2022 merged ``investment banking and securities dealing'' and
``securities brokerage'' into one code, 523150, and ``portfolio management'' and ``investment advice'' into 523940: the
codes that once told a dealer from a broker no longer do. The SEC assigns an SIC code only to issuers of registered
securities, so a private firm on EDGAR has none, however large it is. Codes are the first reading, never the last.

Registrations say more. A broker-dealer files an annual audited report (Form X-17A-5) and periodic financial reports
(FOCUS), of which only the audited statement of financial condition is public; an institutional investment manager with more than \$100 million of US equities files quarterly holdings on
Form 13F; an issuer files annual and quarterly reports; an operator of an alternative trading system files Form ATS-N.
These are obligations, not choices, and they follow the business, not the brand: each is a fact about what the entity
does. The legal entity identifier (Book~2, chapter~28) ties the entities of a group together, and its reference data are
free to reuse.

\begin{method}[Placing an employer from public records]\label{met:in:a-map-of-the-industry:place}
\begin{enumerate}
\item List the group's legal entities: the LEI records, EDGAR's entity search, the national company registries.
\item Read the industry codes, knowing their limits: SIC only for issuers, NAICS coarse, NACE coarser.
\item List each entity's registrations and the families of forms it files: broker-dealer reports, holdings reports,
issuer reports, venue-operator forms.
\item Read the firm's own statement of what it does, on its website and in any filing that describes the business.
\item Decide the business model of the entity the job belongs to, and write down which record decided it.
\end{enumerate}
\end{method}

\begin{example}[One firm, three codes]\label{ex:in:a-map-of-the-industry:codes}
Virtu Financial, the listed holding company of an electronic market maker, carries SIC code 6211, Security Brokers,
Dealers \& Flotation Companies, on EDGAR. Its US operating subsidiary, a broker-dealer, has no SIC code at all but files
broker-dealer annual reports and an alternative-trading-system form. The same SIC code 6211 is carried by Goldman Sachs
and by BlackRock, a bank holding company and an asset manager. The code puts three different employers in one box; the
forms put the market maker's subsidiary in the right one.
\end{example}

\section{Tutorial: classifying employers from their filings}\label{tut:in:a-map-of-the-industry:tut}

\begin{tutorial}
\textbf{Goal.} Measure how well the public codes and the forms an entity files recover its business model. \textbf{End
state:} the table of scores below and \cref{fig:in:a-map-of-the-industry:acc}.
\begin{tutsteps}
\item \textbf{The sample.} \texttt{data/industry/employer\_sample.csv} holds 21 entities that file on EDGAR: eleven
market-making entities, four systematic funds' advisers, two multi-manager platforms, two banks, an asset manager and an
exchange. For each, the SIC code and the families of forms in its recent filing history come from EDGAR's submissions
data; the business model comes from the firm's own statement, with a ledger row per entity.
\item \textbf{The rules.} \texttt{firm.industrymap} implements three rules
(\cref{lst:in:a-map-of-the-industry:rules}): the SIC code alone; the forms alone (a broker-dealer that is not an issuer is
a \omterm{def:in:a-map-of-the-industry:ptf}{principal trading firm}; a filer of holdings reports that is neither is an investment manager); and the forms first,
with SIC for issuers.
\omcode{code/firm/industrymap/firm_industrymap.py}{84}{117}{Three rules that classify an entity from its public
records.}\label{lst:in:a-map-of-the-industry:rules}
\item \textbf{Score.} \texttt{in\_map.accuracy()} scores each rule at two levels: the coarse group (principal trading,
investment manager, bank, infrastructure) and the fine business model.
\item \textbf{Read the misses.} \texttt{in\_map.misses()} lists the entities the best rule gets wrong.
\end{tutsteps}
\end{tutorial}

\begin{center}\small
\begin{tabular}{@{}lrr@{}}
\toprule
rule & right group & right model\\
\midrule
SIC code alone & 3 of 21 & 3 of 21\\
forms filed alone & 15 of 21 & 13 of 21\\
forms, then SIC for issuers & 18 of 21 & 16 of 21\\
\bottomrule
\end{tabular}
\end{center}

\begin{omfigure}
\begin{tikzpicture}
\begin{axis}[width=10.5cm,height=5cm,xbar,area legend,bar width=7pt,xmin=0,xmax=100,xlabel={\small share of the sample classified correctly (\%)},
  ytick={0,1,2},yticklabels={SIC code alone,forms filed alone,{forms, then SIC}},y dir=reverse,ymin=-0.6,ymax=2.6,
  tick label style={font=\footnotesize},grid=major,grid style={gray!20},table/col sep=comma,
  legend style={legend cell align=left,font=\scriptsize,at={(0.5,-0.32)},anchor=north,/tikz/every even column/.append style={column sep=8pt}},legend columns=2]
\addplot[fill=omDef!55,draw=omDef] table[y=k,x=coarse]{figdata/industry/01-a-map-of-the-industry/accuracy.csv};
\addplot[fill=omThm!55,draw=omThm] table[y=k,x=fine]{figdata/industry/01-a-map-of-the-industry/accuracy.csv};
\legend{right group,right business model}
\end{axis}
\end{tikzpicture}

\omcaption{Classifying 21 EDGAR entities from their public records: the SIC code recovers 14\% of them, the forms they
file 71\% of the groups, and the forms with SIC for issuers 86\%. Data: EDGAR submissions and the firms' own
statements, through \texttt{in\_map.accuracy}.}\label{fig:in:a-map-of-the-industry:acc}
\end{omfigure}

The SIC code places only the five issuers, and two of those wrongly: an investment bank and an asset manager sit under
the dealers' code. The forms place every private market maker whose broker-dealer carries the firm's name. The three
misses of the best rule are instructive. One quantitative trading firm's entity in the sample files only holdings
reports: the entity chosen was not a broker-dealer, and the rule saw a fund manager. And the issuer's SIC code misplaces
the bank and the asset manager. Each miss is a place where the method's step~1, listing the group's entities, was
skipped.

\textbf{What to change next.} Add a rule that reads an issuer with SIC 6211 and holdings reports as an asset manager, and
see what it fixes and what it breaks (exercise~7); add ten entities of your own choosing, and ask whether your choice of
entities flatters the rules (exercise~8).

\section{Build: the employer map}\label{bld:in:a-map-of-the-industry:industrymap}

\begin{build}
\bfield{Purpose} One taxonomy of employers for the whole book, so that pay data (chapter~14), firm profiles (chapters
2--10) and roles (Part~III) are sorted the same way.
\bfield{Interface} \texttt{firm.industrymap}: \texttt{MODELS} (business model to group, whose capital, who pays, holding
period), \texttt{GROUP}, \texttt{BOOKS}; \texttt{Entity(name, cik, sic, families, truth)}, \texttt{load\_sample(path)};
\texttt{sic\_class}, \texttt{forms\_class}, and the rules \texttt{by\_sic}, \texttt{by\_forms}, \texttt{combined};
\texttt{score(entities, rule, level)} and \texttt{confusion(entities, rule, level)}.
\bfield{Rules} A rule answers a business model or \texttt{unknown}, never a guess without a record behind it; every
business model has a group and a list of books; a score states its level.
\bfield{Acceptance tests} \texttt{code/firm/industrymap/tests/}: every model has a group and books; SIC and forms rules on
constructed entities; the combined rule falls back to SIC for issuers; scores at the two levels differ where the rule
cannot separate a platform from a systematic fund.
\bfield{Stretch} Add a rule on the NAICS code an employer reports in its labour-condition filings (chapter~14), and on a
UK firm's permissions on the FCA register.
\end{build}

\begin{omsources}
\item US Census Bureau, North American Industry Classification System 2022: codes, index file and 2017-to-2022
concordance.
\item SEC EDGAR submissions data (\texttt{data.sec.gov}) for the 21 entities of the sample; Forms 10-K for 2025 of
JPMorgan Chase, Goldman Sachs, BlackRock, CME Group and Virtu Financial.
\item The firms' own home pages for their descriptions of their business (ledger rows F4--F19).
\item FINRA, Industry Snapshot 2025, tables 2.1.3, 2.1.5 and 2.6.1--2.6.3.
\item FIA, Principal Traders Group and European Principal Traders Association pages; GLEIF, LEI data terms of use.
\end{omsources}

\section{Exercises}

\begin{exercise}[$\star$]\label{exo:in:a-map-of-the-industry:1}
Give the NAICS 2022 code of a firm that makes markets in securities, of one that trades only for its own account, and of
an investment adviser, and say which of them the 2017 codes would have separated.
\end{exercise}

\begin{exercise}[$\star$]\label{exo:in:a-map-of-the-industry:2}
From the dated box, what share of FINRA-registered broker-dealers at the end of 2024 had 150 registered representatives or
fewer, and what share of the proprietary-trading and market-making segments did?
\end{exercise}

\begin{exercise}[$\star$]\label{exo:in:a-map-of-the-industry:3}
Answer the five questions of section~1 for an ETF issuer's index fund, a pension fund's in-house equity team and a
crypto exchange.
\end{exercise}

\begin{exercise}[$\star\star$]\label{exo:in:a-map-of-the-industry:4}
The SIC rule classifies 3 of the 21 entities correctly. How many does it leave as unknown, and why?
\end{exercise}

\begin{exercise}[$\star\star$]\label{exo:in:a-map-of-the-industry:5}
From 2015 to 2024, by how much did the number of FINRA-registered broker-dealers fall, and the number of firms
registered only as investment advisers rise, in per cent?
\end{exercise}

\begin{exercise}[$\star\star$]\label{exo:in:a-map-of-the-industry:6}
In \cref{fig:in:a-map-of-the-industry:acc}, why is the ``right model'' bar lower than the ``right group'' bar for the
forms rule, and which two entities make the difference?
\end{exercise}

\begin{exercise}[$\star\star\star$]\label{exo:in:a-map-of-the-industry:7}
\emph{Coding.} Extend the combined rule so that an issuer with SIC 6211 that also files holdings reports is an asset
manager. How many of the 21 entities does it now place in the right group, and which entity does the new rule still get
wrong, and why?
\end{exercise}

\begin{exercise}[$\star\star\star$]\label{exo:in:a-map-of-the-industry:8}
\emph{Find the flaw.} ``The forms rule places 86\% of employers in the right group, so a job seeker can classify any
trading firm from its EDGAR filings.''
\end{exercise}

\section{Problem: Classifying from the Outside}

\begin{problem}[{Weekend problem --- classifying from the outside}]\label{pb:in:a-map-of-the-industry:1}
A careers adviser wants a tool that tells students what kind of firm has offered them a job, from public records only.
You have the chapter's sample and its rules.

\medskip\noindent\textbf{Part I --- The questions.}
\begin{enumerate}
\item State the five questions that place an employer, and say which three define its business model.
\item Answer them for a market maker and for a systematic fund.
\item Define a \omterm{def:in:a-map-of-the-industry:ptf}{principal trading firm}, and say what the two associations' membership rules add to Book~1's proprietary
trading firm.
\item Why does the chapter classify legal entities or business lines rather than brands?
\item Give the four groups of the taxonomy and the number of business models in each.
\end{enumerate}

\medskip\noindent\textbf{Part II --- The codes.}
\begin{enumerate}[resume]
\item Which NAICS 2022 code covers market making for securities, and which covers speculators for own account?
\item Which 2017 codes were merged into 523150 and into 523940?
\item Why does a private market maker's broker-dealer have no SIC code on EDGAR?
\item Which three entities of the sample carry SIC code 6211, and what are their business models?
\item How many entities of the sample does the SIC rule place correctly, and how many does it leave unknown?
\end{enumerate}

\medskip\noindent\textbf{Part III --- The forms.}
\begin{enumerate}[resume]
\item Name the form families the rules read, and what each proves about the entity.
\item How many entities does the forms rule place in the right group, and in the right business model?
\item Why can the forms rule not tell a multi-manager platform from a systematic fund?
\item How many entities does the combined rule place correctly, at each level?
\item Name the three misses of the combined rule at the group level and the cause of each.
\end{enumerate}

\medskip\noindent\textbf{Part IV --- The verdict.}
\begin{enumerate}[resume]
\item How many broker-dealers did FINRA count at the end of 2024, and how many in its proprietary-trading and
market-making segments?
\item Why is a FINRA size band a poor measure of a market maker's headcount?
\item State the \emph{named result}: the share of the sample each rule classifies correctly, at the group level.
\item Which step of \cref{met:in:a-map-of-the-industry:place} would have caught every miss?
\item In two sentences, what should the adviser's tool report besides its answer?
\end{enumerate}
\end{problem}

\section{Interview questions}

\begin{interviewq}[$\star$ \iqroles{trader, researcher}]\label{iq:in:a-map-of-the-industry:1}
What distinguishes a market maker from a hedge fund? Give three differences that matter to someone who works there.
\end{interviewq}

\begin{interviewq}[$\star$ \iqroles{trader}]\label{iq:in:a-map-of-the-industry:2}
Why does it matter to a trader who pays the firm?
\end{interviewq}

\begin{interviewq}[$\star\star$ \iqroles{developer, risk}]\label{iq:in:a-map-of-the-industry:3}
An entity files annual X-17A-5 reports with the SEC. What does that tell you about it, and what does it not?
\end{interviewq}

\begin{interviewq}[$\star\star$ \iqroles{researcher}]\label{iq:in:a-map-of-the-industry:4}
Estimate how many people work for \omterm{def:in:a-map-of-the-industry:ptf}{principal trading firms} in the United States. State your method and your range.
\end{interviewq}

\begin{interviewq}[$\star\star$ \iqroles{bank, risk}]\label{iq:in:a-map-of-the-industry:5}
A bank's flow desk and a market maker both hold inventory to serve buyers and sellers. Whose capital is at risk in each,
and what limits how much each can hold?
\end{interviewq}

\begin{interviewq}[$\star\star\star$ \iqroles{researcher, mle}]\label{iq:in:a-map-of-the-industry:6}
You build a classifier that labels employers' business models from their public filings and it scores 86\% on a sample
you assembled. How would you know whether to trust the number?
\end{interviewq}
